\documentclass[10pt]{extarticle}
\usepackage[a4paper,margin=1.5cm]{geometry}
\usepackage[T1]{fontenc}
\usepackage{amsmath,amssymb}
\usepackage{graphicx}
\usepackage[export]{adjustbox}
\usepackage{multicol}
\usepackage{enumitem}
\usepackage{booktabs,array,tabularx}
\setlength{\columnsep}{0.7cm}
\setlength{\parindent}{0pt}
\setlength{\parskip}{3pt}
\pagestyle{empty}
\makeatletter
\newcommand{\qfigdef}[4]{\@namedef{qf@file@#1}{#2}\@namedef{qf@w@#1}{#3}\@namedef{qf@wide@#1}{#4}}
\newcommand{\qf@check}[1]{\@ifundefined{qf@file@#1}{\PackageError{score2}{Unknown figure id #1}{}}{}}
\newcommand{\qf@img}[4][]{\includegraphics[width=#2\dimexpr\@nameuse{qf@w@#4}pt\relax,max width=#3,max height=0.45\textheight,#1]{\@nameuse{qf@file@#4}}}
\newif\ifqf@full
\newif\ifqf@boxed
\newsavebox{\qf@box}
\newcommand{\qf@open}{\global\qf@fulltrue\par\global\qf@boxedtrue
  \setbox\qf@box\vbox\bgroup\hsize\textwidth\linewidth\textwidth}
\newcommand{\qf@emit}[1]{\par\egroup\global\qf@boxedfalse
  \ifdim\dimexpr\ht\qf@box+\dp\qf@box\relax<#1\textheight\noindent\usebox{\qf@box}\par\else\unvbox\qf@box\fi}
\newcommand{\qfullsplit}{\ifqf@boxed\qf@emit{0.85}\fi}
\newcommand{\qf@close}{\ifqf@boxed\qf@emit{0.45}\fi\global\qf@fullfalse}
\newcommand{\qfullbegin}{\end{multicols}\qf@open}
\newcommand{\qfullend}{\qf@close\begin{multicols}{2}}
\newcommand{\qfullstart}{\qf@open}
\newcommand{\qfulldone}{\qf@close}
\newcommand{\qfig}[2][1]{\qf@check{#2}%
  \ifqf@full\qf@blk{#1}{#2}\else
  \ifnum\@nameuse{qf@wide@#2}=1 \expandafter\@firstoftwo\else\expandafter\@secondoftwo\fi
  {\end{multicols}\begin{center}\qf@img{#1}{\textwidth}{#2}\end{center}\begin{multicols}{2}}%
  {\qf@blk{#1}{#2}}\fi}
\newcommand{\qf@blk}[2]{\par\smallskip{\centering\qf@img{#1}{\linewidth}{#2}\par}\smallskip}
\newcommand{\qfigi}[2][1]{\qf@check{#2}\qf@img[valign=c]{#1}{\linewidth}{#2}}
\makeatother
\newcommand{\papersec}[1]{\par\bigskip{\centering\Large\bfseries PAPER-#1\par}\medskip}
\newcommand{\partsec}[2]{\par\medskip\noindent{\large\bfseries PART-#1: #2}\par\smallskip}
\newcommand{\sect}[1]{\par\smallskip\noindent\textbf{#1}\par}
\newcommand{\qn}[1]{\par\smallskip\noindent\textbf{#1.}~}
\newcommand{\paraq}[2]{\par\smallskip\noindent\textbf{Paragraph for Questions #1 and #2}\par}
\newenvironment{qoptions}{\begin{itemize}[leftmargin=1.8em,itemsep=1pt,topsep=2pt,label={}]}{\end{itemize}}
\newcommand{\opt}[2]{\item[(#1)] #2}
% figure registry (generated by score2_tools.py finalize)
\qfigdef{p1_phy_s1_q04_1}{Score-II - Test 24 (17 Mar 2024)-img/p1_phy_s1_q04_1.png}{77.8}{0}
\qfigdef{p1_phy_s1_q05_1}{Score-II - Test 24 (17 Mar 2024)-img/p1_phy_s1_q05_1.png}{53.1}{0}
\qfigdef{p1_phy_s1_q06_1}{Score-II - Test 24 (17 Mar 2024)-img/p1_phy_s1_q06_1.png}{119.7}{0}
\qfigdef{p1_phy_s1_q07_1}{Score-II - Test 24 (17 Mar 2024)-img/p1_phy_s1_q07_1.png}{105.5}{0}
\qfigdef{p1_phy_s1_q08_1}{Score-II - Test 24 (17 Mar 2024)-img/p1_phy_s1_q08_1.png}{100.6}{0}
\qfigdef{p1_phy_s1_q10_1}{Score-II - Test 24 (17 Mar 2024)-img/p1_phy_s1_q10_1.png}{98.9}{0}
\qfigdef{p1_phy_s1_q12_1}{Score-II - Test 24 (17 Mar 2024)-img/p1_phy_s1_q12_1.png}{37.2}{0}
\qfigdef{p1_phy_s1_q17_1}{Score-II - Test 24 (17 Mar 2024)-img/p1_phy_s1_q17_1.png}{145.0}{0}
\qfigdef{p1_phy_s1_q19_1}{Score-II - Test 24 (17 Mar 2024)-img/p1_phy_s1_q19_1.png}{85.2}{0}
\qfigdef{p1_phy_s1_q20_1}{Score-II - Test 24 (17 Mar 2024)-img/p1_phy_s1_q20_1.png}{108.3}{0}
\qfigdef{p1_phy_s2_q04_1}{Score-II - Test 24 (17 Mar 2024)-img/p1_phy_s2_q04_1.png}{153.3}{0}
\qfigdef{p1_chem_s1_q13_1}{Score-II - Test 24 (17 Mar 2024)-img/p1_chem_s1_q13_1.png}{55.4}{0}
\qfigdef{p1_chem_s1_q14_1}{Score-II - Test 24 (17 Mar 2024)-img/p1_chem_s1_q14_1.png}{110.4}{0}
\qfigdef{p1_chem_s1_q14_2}{Score-II - Test 24 (17 Mar 2024)-img/p1_chem_s1_q14_2.png}{90.6}{0}
\qfigdef{p1_chem_s1_q14_3}{Score-II - Test 24 (17 Mar 2024)-img/p1_chem_s1_q14_3.png}{67.3}{0}
\qfigdef{p1_chem_s1_q15_1}{Score-II - Test 24 (17 Mar 2024)-img/p1_chem_s1_q15_1.png}{163.6}{0}
\qfigdef{p1_chem_s1_q15_2}{Score-II - Test 24 (17 Mar 2024)-img/p1_chem_s1_q15_2.png}{107.9}{0}
\qfigdef{p1_chem_s1_q15_3}{Score-II - Test 24 (17 Mar 2024)-img/p1_chem_s1_q15_3.png}{136.6}{0}
\qfigdef{p1_chem_s1_q15_4}{Score-II - Test 24 (17 Mar 2024)-img/p1_chem_s1_q15_4.png}{138.4}{0}
\qfigdef{p1_chem_s1_q15_5}{Score-II - Test 24 (17 Mar 2024)-img/p1_chem_s1_q15_5.png}{110.4}{0}
\qfigdef{p1_chem_s1_q16_1}{Score-II - Test 24 (17 Mar 2024)-img/p1_chem_s1_q16_1.png}{85.2}{0}
\qfigdef{p1_chem_s1_q17_1}{Score-II - Test 24 (17 Mar 2024)-img/p1_chem_s1_q17_1.png}{20.8}{0}
\qfigdef{p1_chem_s1_q18_1}{Score-II - Test 24 (17 Mar 2024)-img/p1_chem_s1_q18_1.png}{77.0}{0}
\qfigdef{p1_chem_s1_q18_2}{Score-II - Test 24 (17 Mar 2024)-img/p1_chem_s1_q18_2.png}{105.2}{0}
\qfigdef{p1_chem_s1_q18_3}{Score-II - Test 24 (17 Mar 2024)-img/p1_chem_s1_q18_3.png}{89.8}{0}
\qfigdef{p1_chem_s1_q18_4}{Score-II - Test 24 (17 Mar 2024)-img/p1_chem_s1_q18_4.png}{115.5}{0}
\qfigdef{p1_chem_s1_q19_1}{Score-II - Test 24 (17 Mar 2024)-img/p1_chem_s1_q19_1.png}{84.6}{0}
\qfigdef{p1_chem_s1_q19_2}{Score-II - Test 24 (17 Mar 2024)-img/p1_chem_s1_q19_2.png}{19.1}{0}
\qfigdef{p1_chem_s1_q19_3}{Score-II - Test 24 (17 Mar 2024)-img/p1_chem_s1_q19_3.png}{90.8}{0}
\qfigdef{p1_chem_s1_q19_4}{Score-II - Test 24 (17 Mar 2024)-img/p1_chem_s1_q19_4.png}{19.8}{0}
\qfigdef{p1_chem_s2_q07_1}{Score-II - Test 24 (17 Mar 2024)-img/p1_chem_s2_q07_1.png}{24.0}{0}
\qfigdef{p1_chem_s2_q08_1}{Score-II - Test 24 (17 Mar 2024)-img/p1_chem_s2_q08_1.png}{161.1}{0}
\makeatletter\@beginparpenalty=10000\makeatother
\begin{document}
\begin{center}{\Large\bfseries JEE (Main+Advanced) Score-II --- Test 24 (17 Mar 2024)}\end{center}
\begin{multicols}{2}
\raggedcolumns
\partsec{1}{PHYSICS}
\sect{Section-I}

\qn{1} The gravitational field in a region is given by $\vec E=-(60\ \mathrm{N/kg})(\hat i+\hat j)$. Find the work done in shifting a particle of mass 3 kg from origin (0, 0) to a point (3, 7).
\begin{qoptions}
\opt{A}{$-1800$ J}
\opt{B}{1800 J}
\opt{C}{$-900$ J}
\opt{D}{900 J}
\end{qoptions}

\qn{2} Equal masses of three liquids A, B and C have temperatures $10^\circ$C, $25^\circ$C and $40^\circ$C respectively. If A and B are mixed, the mixture has a temperature of $15^\circ$C. If B and C are mixed, the mixture has a temperature of $30^\circ$C. If A and C are mixed, the mixture will have a temperature of :-
\begin{qoptions}
\opt{A}{$16^\circ$C}
\opt{B}{$20^\circ$C}
\opt{C}{$25^\circ$C}
\opt{D}{$29^\circ$C}
\end{qoptions}

\qn{3} \textbf{Assertion (A) :} Under electrostatic condition, the electric field at surface of conductor can have no component parallel to surface.\\
\textbf{Reason (R) :} The surface of conductor is equipotential and no work is done by electric field to move the charge on surface.
\begin{qoptions}
\opt{A}{(A) is True, (R) is True ; (R) is a correct explanation for (A).}
\opt{B}{(A) is True, (R) is True ; (R) is not a correct explanation for (A).}
\opt{C}{(A) is True, (R) is False.}
\opt{D}{(A) is False, (R) is True.}
\end{qoptions}

\qn{4} A thermodynamic system undergoes cyclic process ABCDA as shown in figure. The work done by the system is :-
\qfig{p1_phy_s1_q04_1}
\begin{qoptions}
\opt{A}{$P_0V_0$}
\opt{B}{$2P_0V_0$}
\opt{C}{$\dfrac{P_0V_0}{2}$}
\opt{D}{Zero}
\end{qoptions}

\qn{5} A solid conducting sphere of radius $a$ has a net positive charge of $2Q$. A conducting spherical shell of inner radius $b$ and outer radius $c$ is concentric with the solid sphere and has a net charge of $-Q$. The surface charge density on the inner and outer surfaces of the spherical shell will be :-
\qfig{p1_phy_s1_q05_1}
\begin{qoptions}
\opt{A}{$-\dfrac{2Q}{4\pi b^2},\ \dfrac{Q}{4\pi c^2}$}
\opt{B}{$-\dfrac{Q}{4\pi b^2},\ \dfrac{Q}{4\pi c^2}$}
\opt{C}{$0,\ \dfrac{Q}{4\pi c^2}$}
\opt{D}{$\dfrac{2Q}{4\pi b^2},\ -\dfrac{Q}{4\pi c^2}$}
\end{qoptions}

\qn{6} The figure shows a circuit with E as the earthing of the common plate. The potentials at P and Q are :-
\qfig{p1_phy_s1_q06_1}
\begin{qoptions}
\opt{A}{0 V, $-1000$ V}
\opt{B}{1000 V, 0 V}
\opt{C}{$+600$ V, $-400$ V}
\opt{D}{$+400$ V, $-600$ V}
\end{qoptions}

\qn{7} A frictionless cart A of mass 100 kg carries other two frictionless carts B and C having masses 8 kg and 4 kg respectively connected by a string passing over a pulley as shown in the figure. What horizontal force $F$ must be applied on the cart A so that smaller carts do not move relative to it ? (Take : $g=10\ \mathrm{m/s^2}$)
\qfig{p1_phy_s1_q07_1}
\begin{qoptions}
\opt{A}{150 N}
\opt{B}{340 N}
\opt{C}{560 N}
\opt{D}{630 N}
\end{qoptions}

\qn{8} An irregular loop of conducting wire is lying on a frictionless table as shown in the figure. The wire is clamped at points a and k, when a current $i$ is passed through it, then the wire will :-
\qfig{p1_phy_s1_q08_1}
\begin{qoptions}
\opt{A}{Become parallel in two parts}
\opt{B}{Collapse more}
\opt{C}{Form a circular loop}
\opt{D}{No change in the configuration}
\end{qoptions}

\qn{9} A capacitor of capacitance 100 $\mu$F is charged to potential difference of 24 V. The charging battery is disconnected and the capacitor is connected to another battery of emf 12 V with positive plate of capacitor joined with positive terminal of battery.
\par\nobreak
\begin{center}
\begin{tabular}{|c|p{4.6cm}|c|c|}\hline
 & \multicolumn{1}{c|}{\textbf{List--I}} & & \textbf{List--II}\\\hline
(P) & Work done by battery of emf 24 V & (1) & 7.2 mJ\\\hline
(Q) & The decrease in electrostatic field energy & (2) & 14.4 mJ\\\hline
(R) & Work done on battery of emf 12 V & (3) & 21.6 mJ\\\hline
(S) & The heat developed during flow of charge after reconnection. & (4) & 57.6 mJ\\\hline
\end{tabular}
\end{center}
\begin{qoptions}
\opt{A}{(P) - (4) ; (Q) - (3) ; (R) - (2) ; (S) - (1)}
\opt{B}{(P) - (4) ; (Q) - (2) ; (R) - (1) ; (S) - (3)}
\opt{C}{(P) - (2) ; (Q) - (3) ; (R) - (1) ; (S) - (4)}
\opt{D}{(P) - (2) ; (Q) - (3) ; (R) - (4) ; (S) - (1)}
\end{qoptions}

\qn{10} In the adjoining circuit, initially the switch S is open. The switch `S' is closed at $t=0$. The difference between the maximum and minimum current that can flow through the battery is :-
\qfig{p1_phy_s1_q10_1}
\begin{qoptions}
\opt{A}{2 A}
\opt{B}{3 A}
\opt{C}{1 A}
\opt{D}{0.5 A}
\end{qoptions}

\qn{11} A vernier callipers has a least count of 1mm of its main scale. If 11 divisions of its main scale coincides with 12 divisions of vernier scale then its least count will be :-
\begin{qoptions}
\opt{A}{$\dfrac{1}{12}$ mm}
\opt{B}{$\dfrac{1}{11}$ mm}
\opt{C}{$\dfrac{11}{12}$ mm}
\opt{D}{$\dfrac{12}{11}$ mm}
\end{qoptions}

\qn{12} A string is suspended from a rigid support as shown in the figure. Two impulses are given at ends A \& B one by one, which results travelling transverse pulses from A to B \& B to A. If the time taken by the pulses travelling from A to B is $T_1$ \& by the pulse travelling from B to A is $T_2$ then :-
\qfig{p1_phy_s1_q12_1}
\begin{qoptions}
\opt{A}{$T_1=T_2$}
\opt{B}{$T_1>T_2$}
\opt{C}{$T_2>T_1$}
\opt{D}{$T_1$ may be either less than or greater than $T_2$.}
\end{qoptions}

\qn{13} Calculate the binding energy of an alpha particle (in MeV) (approx) from the following data :-\\
mass of H atom = 1.007862 u\\
mass of neutron = 1.008656 u\\
mass of He atom = 4.00206 u\\
Take 1 u = 931 MeV/c$^2$.
\begin{qoptions}
\opt{A}{25}
\opt{B}{40}
\opt{C}{29}
\opt{D}{19}
\end{qoptions}

\qn{14} A lift is raised by light cables from rest to rest through a height 180 m. The greatest tension which the cables can safely bear is 9 times of lift's weight, if ascent is to made in shortest time then maximum speed of lift during its journey will be :- ($g=10\ \mathrm{m/s^2}$)
\begin{qoptions}
\opt{A}{$40\sqrt2$ m/s}
\opt{B}{$30\sqrt5$ m/s}
\opt{C}{$25\sqrt7$ m/s}
\opt{D}{$45\sqrt3$ m/s}
\end{qoptions}

\qn{15} \textbf{Statement-1 :} The mass of nucleus is not integral multiples of the mass of single nucleon.\\
\textbf{Statement-2 :} Due to mass defect the total mass of nucleus is not equal to masses of its constituents.
\begin{qoptions}
\opt{A}{Statement-1 is True, Statement-2 is True ; Statement-2 is a correct explanation for Statement-1.}
\opt{B}{Statement-1 is True, Statement-2 is True ; Statement-2 is NOT a correct explanation for Statement-1.}
\opt{C}{Statement-1 is True, Statement-2 is False.}
\opt{D}{Statement-1 is False, Statement-2 is True.}
\end{qoptions}

\qn{16} A point object is placed on optical axis of a convex lens of focal length 10 cm at a distance 20 cm to left of it. An eye is placed at a distance 30 cm to right of lens and a distance `$y$' below the optical axis. If the diameter of lens is 4 cm then maximum value of `$y$' to see the image is :- (assume parallel approximation)
\begin{qoptions}
\opt{A}{4 cm}
\opt{B}{2 cm}
\opt{C}{1 cm}
\opt{D}{$\dfrac43$ cm}
\end{qoptions}

\qn{17} A part of circuit is shown in figure. As the switch S is closed in the circuit, the magnitude of change in potential of point `P' will be :-
\qfig{p1_phy_s1_q17_1}
\begin{qoptions}
\opt{A}{11 V}
\opt{B}{9 V}
\opt{C}{7 V}
\opt{D}{6 V}
\end{qoptions}

\qn{18} Mark out the correct option :-
\begin{qoptions}
\opt{A}{The energies of all the small parts of equal length are equal in travelling wave.}
\opt{B}{The energies of all the small parts of equal length are equal in standing wave.}
\opt{C}{The phase difference between two particles may be $\dfrac{\pi}{2}$ in travelling wave.}
\opt{D}{The phase difference between two particles may be $\dfrac{\pi}{2}$ in standing wave.}
\end{qoptions}

\qn{19} The plank in figure moves a distance 100 mm to right while the centre of mass of sphere of radius 150 mm moves 75 mm to the left. If there is no slipping anywhere then angular displacement of sphere (in radian) is
\qfig{p1_phy_s1_q19_1}
\begin{qoptions}
\opt{A}{$\dfrac76$}
\opt{B}{1}
\opt{C}{$\dfrac12$}
\opt{D}{$\dfrac16$}
\end{qoptions}

\qn{20} The truth table for digital circuit shown in figure is :-
\qfig{p1_phy_s1_q20_1}
\begin{qoptions}
\opt{A}{\begin{tabular}[t]{|c|c|c|}\hline A&B&Y\\\hline 0&0&0\\0&1&0\\1&0&0\\1&1&1\\\hline\end{tabular}}
\opt{B}{\begin{tabular}[t]{|c|c|c|}\hline A&B&Y\\\hline 0&0&1\\0&1&0\\1&0&0\\1&1&0\\\hline\end{tabular}}
\opt{C}{\begin{tabular}[t]{|c|c|c|}\hline A&B&Y\\\hline 0&0&1\\0&1&1\\1&0&1\\1&1&0\\\hline\end{tabular}}
\opt{D}{\begin{tabular}[t]{|c|c|c|}\hline A&B&Y\\\hline 0&0&0\\0&1&1\\1&0&1\\1&1&0\\\hline\end{tabular}}
\end{qoptions}

\sect{Section-II}

\qn{1} A boy of mass 50 kg stands on a board of mass 150 kg which in turn is kept on a frictionless horizontal ice surface. The boy makes a jump with a velocity component 15 m/s in horizontal direction with respect to the ice. With what rate in m/s are the boy and the board separating from each other?

\qn{2} Current in an A.C. circuit is given by $i=60\pi\sin(\pi t+\pi/4)$, then the average value of current during time $t=0$ to $t=3$ sec is :-

\qn{3} A small particle of mass $m$ moves in such a way that the potential energy $U=ar^3$ where $a$ is a constant and $r$ is the distance of the particle from the origin. Assuming Bohr's model of quantization of angular momentum and circular orbits, if the radius of $n^{\text{th}}$ allowed orbit is proportional to $n^x$. Then the value of $100x$ is :-

\qn{4} Water flows in a horizontal tube as shown in figure. The pressure of water changes by 1500 N/m$^2$ between A and B where the areas of cross-section are 60 cm$^2$ and 30 cm$^2$ respectively. Find the rate of flow of water through the tube in cm$^3$/s.
\qfig{p1_phy_s2_q04_1}

\qn{5} Find out the moment of inertia of a semi circular disc of mass 10 kg and radius 3m about an axis passing through its centre of mass and perpendicular to the plane? (Take $\pi^2=10$)

\qn{6} A paramagnetic material has $10^{20}$ atoms/m$^3$. It's magnetic susceptibility at temperature 200 k is $1.2\times10^{-4}$. It's susceptibility at 240 K is $x$ then $10^4x$ will be :-

\qn{7} A transverse wave $y=0.021\sin\left(\dfrac{x}{20}+40t\right)$ where $x$ and $y$ are in meter and $t$ is in second is propagating along a wire of area of cross-section $5\times10^{-4}$ cm$^2$ and density $2\times10^3$ kg/m$^3$. Find the tension (in N) in the wire.

\qn{8} An astronomical Telescope has angular magnification of magnitude 5x for different objects. When final image is formed at infinity, its tube length is 36 cm. Find focal length (in cm) of its objective.

\qn{9} A parallel beam of light is incident on a thin film of air of uniform thickness. Successive wavelengths of 7200 \AA\ and 5400 \AA\ are observed to be missing from the spectrum of reflected light normally. Find the other wavelength (in \AA) in the visible region (4000 \AA\ to 7800 \AA) missing in the reflected spectrum.

\qn{10} The longest wavelength of Lyman series for hydrogen atom is the same as the wavelength of certain line in spectrum of He$^+$, when the electron makes a translation from $n\to2$. What is the value of $n$ ?

\partsec{2}{CHEMISTRY}
\sect{Section-I}

\qn{1} Which of the following does not form black precipitate with ammonical solution of H$_2$S.
\begin{qoptions}
\opt{A}{NiCl$_2$}
\opt{B}{ZnSO$_4$}
\opt{C}{FeCl$_2$}
\opt{D}{CoSO$_4$}
\end{qoptions}

\qn{2} Which of the following order is correct?
\begin{qoptions}
\opt{A}{B $<$ Al $<$ Ga $<$ In $<$ T$\ell$ (Atomic radius)}
\opt{B}{B$^{+3}$ $<$ Al$^{+3}$ $<$ Ga$^{+3}$ $<$ In$^{+3}$ $<$ Tl$^{+3}$ (Ionic radius)}
\opt{C}{B $>$ Ga $>$ Al $>$ In $>$ T$\ell$ (Ionization energy)}
\opt{D}{B $>$ Al $>$ Ga $>$ In $>$ T$\ell$ (Electronegativity)}
\end{qoptions}

\qn{3} When native ore of gold treated with aqueous solution of NaCN in the presence of oxygen then which of the following property is not related with the formed complex.
\begin{qoptions}
\opt{A}{sp$^3$ hybridized}
\opt{B}{Diamagnetic}
\opt{C}{Gold is in +1 oxidation state}
\opt{D}{Maximum number of atom in one plane is 5}
\end{qoptions}

\qn{4} Which of the following reagent can not be used for the distinction between Fe$^{+2}$ \& Fe$^{+3}$.
\begin{qoptions}
\opt{A}{Addition of KSCN}
\opt{B}{Addition of NH$_4$OH}
\opt{C}{Addition of dil. HCl}
\opt{D}{Addition of K$_4$[Fe(CN)$_6$]}
\end{qoptions}

\qn{5} \textbf{Assertion (A) :} Mn$^{+3}$ in aqueous solution act as oxidizing agent.\\
\textbf{Reason (R) :} The $\mathrm{E^{o}_{Mn^{+2}/Mn^{+3}}}$ value is negative
\begin{qoptions}
\opt{A}{A is false but R is true.}
\opt{B}{A is true but R is false.}
\opt{C}{Both A and R are true and R is the correct explanation of A.}
\opt{D}{Both A and R are true and R is NOT the correct explanation of A.}
\end{qoptions}

\qn{6} What volume of 0.2N KMnO$_4$ solution is needed for complete reaction with 6.64 g Fe$_{0.9}$O$_{1.0}$, in acidic medium ? (Fe = 56)
\begin{qoptions}
\opt{A}{280 ml}
\opt{B}{$\dfrac{280}{9}$ ml}
\opt{C}{$\dfrac{350}{9}$ ml}
\opt{D}{350 ml}
\end{qoptions}

\qn{7} Carbon monoxide in water gas reacts with steam according to the following reaction:\\
$\mathrm{CO(g)+H_2O(g)\rightleftharpoons CO_2(g)+H_2(g)}$ ; $K_{eq.}=\dfrac13$\\
If two volume of water gas (containing 1 : 1 of CO and H$_2$) are mixed with five volumes of steam, then the volume ratio of the four gases (CO : H$_2$O : CO$_2$ : H$_2$) present at equilibrium is :
\begin{qoptions}
\opt{A}{1 : 9 : 1 : 3}
\opt{B}{1 : 15 : 1 : 5}
\opt{C}{1 : 1 : 1 : 1}
\opt{D}{1 : 3 : 1 : 1}
\end{qoptions}

\qn{8} The heat of ionisation of formic acid is 1.5 kJ/mol. 9.2 g formic acid on reaction with 14 g ammonium hydroxide evolves 10.8 kJ of heat. The heat of ionisation of ammonium hydroxide is :\\
(Heat of neutralisation of strong acid \& strong base $=-57$ kJ/eq.)
\begin{qoptions}
\opt{A}{1.5 kJ/mol}
\opt{B}{2.5 kJ/mol}
\opt{C}{3 kJ/mol}
\opt{D}{10.8 kJ/mol}
\end{qoptions}

\qn{9} For a 1-2 type sparingly soluble electrolyte, solubility product is $4\times10^{-12}$. The concentration of cation in saturated solution will be :
\begin{qoptions}
\opt{A}{$2\times10^{-6}$ M}
\opt{B}{$10^{-4}$ M}
\opt{C}{$2\times10^{-4}$ M}
\opt{D}{$10^{-6}$ M}
\end{qoptions}

\qn{10} Given below are two statements :\\
\textbf{Statement-I :} A mixture of phenol and aniline will show positive deviation from Raoults law .\\
\textbf{Statement-II :} When intermolecular intractions weaken after mixing liquids then vapour pressure increases.\\
In the light of the above statements, choose the most appropriate answer from the options given below.
\begin{qoptions}
\opt{A}{Statement-I is false but Statement-II is true}
\opt{B}{Both Statement-I and Statement-II are false}
\opt{C}{Statement-I is true but Statement-II is false}
\opt{D}{Both Statement-I and Statement-II are true}
\end{qoptions}

\qn{11} Select incorrect statements :
\begin{qoptions}
\opt{A}{Cell potential of mercury cell remain same during its life}
\opt{B}{Cell potential of Ni-Cad cell remain same during its life}
\opt{C}{In fuel cell, oxygen is reduced at cathode}
\opt{D}{In corrosion, $\mathrm{H^{+}_{(aq.)}}$ is reduced at cathode}
\end{qoptions}

\qn{12} Select the incorrect statements :
\begin{qoptions}
\opt{A}{Temperature is the measure of average chaotic motion of particles in system.}
\opt{B}{Heat added to system at higher temperature causes greater randomness.}
\opt{C}{For strongly endothermic reactions , it is unlikely to form much product.}
\opt{D}{Spontaneity of a reaction may change by changing temperature.}
\end{qoptions}

\qn{13} Find out the correct IUPAC name of following compound.
\qfig{p1_chem_s1_q13_1}
\begin{qoptions}
\opt{A}{2-Amino-1-chloro-1-methyl cyclohex-4-ene}
\opt{B}{2-Chloro-2-methyl cyclo hex-4-en-1-amine}
\opt{C}{6-Chloro-6-methyl cyclo hex-3-en-1-amine}
\opt{D}{4-Chloro-4-methyl cyclo hex-1-en-5-amine}
\end{qoptions}

\qn{14} Which of the following gives meso compound as product?
\begin{qoptions}
\opt{A}{$\mathrm{CH_3-CH=CH_2}\ \xrightarrow[\mathrm{NaHSO_3}]{\mathrm{OsO_4}}$}
\opt{B}{\qfigi{p1_chem_s1_q14_1}}
\opt{C}{\qfigi{p1_chem_s1_q14_2}}
\opt{D}{\qfigi{p1_chem_s1_q14_3}}
\end{qoptions}

\qn{15}
\qfig{p1_chem_s1_q15_1}
E and F respectively are
\begin{qoptions}
\opt{A}{\qfigi{p1_chem_s1_q15_2}}
\opt{B}{\qfigi{p1_chem_s1_q15_3}}
\opt{C}{\qfigi{p1_chem_s1_q15_4}}
\opt{D}{\qfigi{p1_chem_s1_q15_5}}
\end{qoptions}

\qn{16}
\qfig{p1_chem_s1_q16_1}
Which statement is INCORRECT ?
\begin{qoptions}
\opt{A}{Reduced product of P and Q will be metamers to each other. (Consider LiAlH$_4$ as a reducing agent)}
\opt{B}{By dry distillation of hydrolysed product of P with Ca(OH)$_2$/$\Delta$ gives benzophenone.}
\opt{C}{One of the hydrolysis product of Q reacts with NaNO$_2$ + HCl followed by reaction with phenol in mild basic medium gives orange red dye.}
\opt{D}{Electrophile involved in the formation of Q is dichlorocarbene.}
\end{qoptions}

\qn{17} In chromatography, which of the following statements is incorrect for R$_f$(Retardation factor) ?
\begin{qoptions}
\opt{A}{R$_f$ value depends on the type of chromatography.}
\opt{B}{Higher R$_f$ value means higher adsorption.}
\opt{C}{R$_f$ value is dependent on the mobile phase.}
\opt{D}{The value of R$_f$ can not be more than one.}
\end{qoptions}

\qn{18} \qfigi{p1_chem_s1_q17_1}\ $\xrightarrow[\mathrm{(ii)\ Zn(CH_3)_2S}]{\mathrm{(i)\ O_3}}$ (A) $\xrightarrow{\ \overset{\ominus}{\mathrm{O}}\mathrm{H},\ \Delta\ }$ (B)\\
Major product A and B are respectively.
\begin{qoptions}
\opt{A}{\qfigi{p1_chem_s1_q18_1}}
\opt{B}{\qfigi{p1_chem_s1_q18_2}}
\opt{C}{\qfigi{p1_chem_s1_q18_3}}
\opt{D}{\qfigi{p1_chem_s1_q18_4}}
\end{qoptions}

\qn{19}
\par\nobreak
\begin{center}
\begin{tabular}{|c|p{3.4cm}|c|p{2.0cm}|}\hline
\multicolumn{2}{|c|}{\shortstack{\textbf{List-I}\\(Reaction)}} & \multicolumn{2}{c|}{\shortstack{\textbf{List-II}\\(Name of reaction)}}\\\hline
(P) & \qfigi{p1_chem_s1_q19_1} & (1) & Gabriel pthalamide synthesis to prepare 1$^\circ$ amine\\\hline
(Q) & \qfigi{p1_chem_s1_q19_2}\newline{\scriptsize $\xrightarrow[\text{(ii) H}^{\oplus}]{\text{(i) CO}_2,\ \text{NaOH},\ \Delta}$} & (2) & Reimer Tiemann reaction\\\hline
(R) & \qfigi{p1_chem_s1_q19_3} & (3) & Kolbe Schmitt reaction\\\hline
(S) & \qfigi{p1_chem_s1_q19_4}\newline{\scriptsize $\xrightarrow[\text{(ii) H}^{\oplus}]{\text{(i) CHCl}_3,\text{NaOH},\Delta}$} & (4) & Cannizaro reaction\\\hline
\end{tabular}
\end{center}
\begin{qoptions}
\opt{A}{P $\to$ 4, Q $\to$ 3, R $\to$ 2, S $\to$ 1}
\opt{B}{P $\to$ 4, Q $\to$ 2, R $\to$ 1, S $\to$ 3}
\opt{C}{P $\to$ 3, Q $\to$ 4, R $\to$ 2, S $\to$ 1}
\opt{D}{P $\to$ 4, Q $\to$ 3, R $\to$ 1, S $\to$ 2}
\end{qoptions}

\qn{20} Hybridisation of (underlined atom) in given species is respectively :\\
(I) $[\underline{\mathrm{Fe}}(\mathrm{CO})_5]$\\
(II) $[\underline{\mathrm{Ni}}(\mathrm{CO})_4]$\\
(III) $\underline{\mathrm{Br}}\mathrm{F}_5$\\
(IV) $\underline{\mathrm{N}}(\mathrm{SiH}_3)_3$
\begin{qoptions}
\opt{A}{dsp$^3$ , dsp$^2$ , sp$^3$d , sp$^3$}
\opt{B}{sp$^3$d , sp$^3$ , sp$^3$d$^2$ , sp$^2$}
\opt{C}{dsp$^3$ , sp$^3$ , sp$^3$d$^2$ , sp$^2$}
\opt{D}{dsp$^3$ , dsp$^2$ , sp$^3$d$^2$ , sp$^3$}
\end{qoptions}

\sect{Section-II}

\qn{1} What is the sum of oxidation state, and Co-ordination number when CoCl$_2$ treated with excess of KNO$_2$.

\qn{2} Find the sum of A--A linkages in the following. (where A is `O' or `S' or `N' or `P')\\
H$_2$O$_2$ , H$_2$S$_2$O$_5$ , N$_2$O$_5$ , P$_4$O$_{10}$ , P$_4$ , HNO$_4$ , H$_2$S$_2$O$_8$

\qn{3} Find the number of chemical species which has bond order equal or greater than 2 and paramagnetic\\
O$_2^+$ , O$_2$ , N$_2^+$ , N$_2^-$ , CN$^-$ , OF , NO$^+$ , C$_2$

\qn{4} The energy (in KJ) required to excite 1L of Hydrogen gas at 1atm and 300K to the first excited state of atomic hydrogen is :\\
(The energy for the dissociation of H--H Bond is 432 KJ mol$^{-1}$ and the ionization energy of Hydrogen is 1296KJ mol$^{-1}$).\\
R = 0.08 L atm/K-mole

\qn{5} For a particular reaction A $\to$ product, graph of reciprocal of concentration of reactant (y-axis) versus time (x-axis) is a straight line having intercept 10 (on y-axis) and slope 0.1. Calculate time (in min) required for 80\% completion of reaction. [Time - min, concentration = molar]

\qn{6} If the pKa of lactic acid is 5, then the pH of 0.005 M calcium lactate solution at 25$^\circ$ C is \underline{\hspace{2em}} $\times10^{-1}$.

\qn{7} How many total allylic monobromo substituted products are formed with following compound when it react with N-bromo succinamide\\
\qfigi{p1_chem_s2_q07_1}\ $\xrightarrow[\Delta]{\mathrm{NBS}}$ Allylic mono bromo substituted products

\qn{8}
\qfig{p1_chem_s2_q08_1}
Molecular weight of final organic product is?

\qn{9} How many statement(s) is/are correct about nucleic acids.\\
(i) Synthesis of nucleic acid is from 5$'$ end to 3$'$ end during replication or transcription.\\
(ii) DNA contains adenine, thymine, guanine and cytosine nitrogenous base and ribose pentose sugar.\\
(iii) RNA contains adenine uracil, guanine and cytosine nitrogenous base and ribose pentose sugar.\\
(iv) Secondary structure of DNA is double Helix structure which has hydrogen bonds.\\
(v) RNA exist as single strand.\\
(vi) r-RNA is the structural part of ribosome.\\
(vii) $\alpha$-helix and $\beta$-pleated sheet are secondary structure of DNA.

\qn{10} Find the number of (atoms/ions) which have exactly half filled d or f subshell\\
Gd$^{+2}$ , Gd$^{+3}$ , Gd , Fe$^{+3}$ , Eu$^{+2}$ , Cr , Yb$^{+2}$ , Tb$^{+4}$ , Mn , Cm$^{+3}$\\
(Atomic Number : Gd = 64, Fe = 26, Eu = 63, Cr = 24, Yb = 70, Tb = 65, Mn = 25, Cm = 96)

\partsec{3}{MATHEMATICS}
\sect{Section-I}

\qn{1} Let $e$ be the eccentricity of a hyperbola and $f(e)$ be the eccentricity of its conjugate hyperbola then $\displaystyle\int\left(f(e)+f(f(e))\right)de=g(e)$ and $g\left(\sqrt2\right)=2$, then $g(e)=$
\begin{qoptions}
\opt{A}{$\dfrac12\sqrt{e^2-1}+\dfrac{e^2}{2}$}
\opt{B}{$\dfrac12\sqrt{e^2+1}+\dfrac{e^2}{2}$}
\opt{C}{$\sqrt{e^2-1}+\dfrac{e^2}{2}$}
\opt{D}{$\sqrt{e^2+1}+\dfrac{e^2}{2}$}
\end{qoptions}

\qn{2} The integral $\displaystyle\int_0^1\dfrac{\tan^{-1}x}{1+x}\,dx$ equals
\begin{qoptions}
\opt{A}{$\pi\log_e2$}
\opt{B}{$\dfrac{\pi}{4}\log_e2$}
\opt{C}{$\dfrac{\pi}{8}\log_e2$}
\opt{D}{$\dfrac{\pi}{2}\log_e2$}
\end{qoptions}

\qn{3} \textbf{Assertion (A) :} Line $5x+7y+9=0$ cannot intersect the parabola $y=x^2$ in a rational point (whose abscissa \& ordinate both are rational)\\
\textbf{Reason (R) :} For odd integer $a,b$ and $c$ equation $ax^2+bx+c=0$ cannot have rational roots.
\begin{qoptions}
\opt{A}{Both \textbf{A} and \textbf{R} are correct and \textbf{R} is the correct explanation of \textbf{A}}
\opt{B}{Both \textbf{A} and \textbf{R} are correct but \textbf{R} is \textbf{NOT} the correct explanation of \textbf{A}}
\opt{C}{\textbf{A} is correct but \textbf{R} is not correct}
\opt{D}{\textbf{A} is not correct but \textbf{R} is correct}
\end{qoptions}

\qn{4} $\displaystyle\int_0^{\pi}x\left(\sin^2(\sin x)+\cos^2(\cos x)\right)dx=$
\begin{qoptions}
\opt{A}{$\pi^2$}
\opt{B}{$\dfrac{\pi^2}{2}$}
\opt{C}{$\dfrac{\pi^2}{4}$}
\opt{D}{$\dfrac{\pi^2}{8}$}
\end{qoptions}

\qn{5} Let $z_1$ and $z_2$ be two complex number such that $|z_1|=1$ and $|z_2|=10$. If $\theta=\arg\left(\dfrac{z_1-z_2}{z_2}\right)$ then maximum value of $\tan^2\theta$ is
\begin{qoptions}
\opt{A}{$\dfrac1{10}$}
\opt{B}{$\dfrac1{100}$}
\opt{C}{$\dfrac1{99}$}
\opt{D}{$\dfrac{10}{99}$}
\end{qoptions}

\qn{6} The probability that the birthdays of six different persons will fall in exactly two calendar months
\begin{qoptions}
\opt{A}{$\dfrac{342}{12^5}$}
\opt{B}{$\dfrac{341}{12^6}$}
\opt{C}{$\dfrac{341}{12^5}$}
\opt{D}{$\dfrac{342}{12^6}$}
\end{qoptions}

\qn{7} The system of equations\\
$(1-\lambda)x+3y-4z=0$\\
$x-(3+\lambda)y+5z=0$\\
$3x+y-\lambda z=0$\\
possesses non-trivial solutions for
\begin{qoptions}
\opt{A}{exactly three distinct real values of $\lambda$}
\opt{B}{exactly two distinct real values of $\lambda$}
\opt{C}{exactly one real values of $\lambda$}
\opt{D}{infinite number of distinct real values of $\lambda$}
\end{qoptions}

\qn{8} If $M=\begin{bmatrix}1&c&3\\1&3&3\\2&4&4\end{bmatrix}$ is the adjoint of a $3\times3$ matrix N and det.(N) $=4$, then $c$ is equal to
\begin{qoptions}
\opt{A}{0}
\opt{B}{4}
\opt{C}{5}
\opt{D}{11}
\end{qoptions}

\qn{9} \par\nobreak
\begin{center}\small
\begin{tabular}{|c|p{3.6cm}|c|c|}\hline
 & \multicolumn{1}{c|}{\textbf{List-I}} & & \textbf{List-II}\\\hline
(A) & If the line $\dfrac{x-4}{1}=\dfrac{y-2}{1}=\dfrac{z-2k}{2}$ lies in the plane $2x-4y+z=3$, then $k$ is equal to & (I) & $\dfrac92$\\\hline
(B) & If the lines $\dfrac{x-1}{2}=\dfrac{y+1}{3}=\dfrac{z-1}{4}$ and $\dfrac{x-3}{1}=\dfrac{y-k}{2}=\dfrac{z}{1}$ intersect, then value of $k$ is & (II) & $\dfrac32$\\\hline
(C) & If a plane passes through the point $(1,1,1)$ cuts the coordinate axes at A, B \& C respectively such that $OA=OB=OC$ where O is the origin, then volume of tetrahedron OABC is equal to & (III) & $\dfrac52$\\\hline
(D) & If a plane P passes through $(1,-2,1)$ is perpendicular to two planes $2x-2y+z=0$ and $x-y+2z=4$. The distance of plane P from the point $\left(-1,\dfrac{5}{\sqrt2},\dfrac{3}{\sqrt2}\right)$ is & (IV) & $\dfrac72$\\\hline
\end{tabular}
\end{center}
\begin{qoptions}
\opt{A}{(A) $-$ II, (B) $-$ I, (C) $-$ I, (D) $-$ III}
\opt{B}{(A) $-$ I, (B) $-$ III, (C) $-$ IV, (D) $-$ II}
\opt{C}{(A) $-$ II, (B) $-$ II, (C) $-$ I, (D) $-$ III}
\opt{D}{(A) $-$ II, (B) $-$ I, (C) $-$ III, (D) $-$ III}
\end{qoptions}

\qn{10} If $\displaystyle\sum_{m=1}^{12}12m\cdot{}^{12}C_m\cdot{}^{11}C_{m-1}$ is equal to $12^p\cdot{}^{22}C_q$, then $(p+q)$ has the value equal to
\begin{qoptions}
\opt{A}{11}
\opt{B}{12}
\opt{C}{13}
\opt{D}{14}
\end{qoptions}

\qn{11} If $f(1)=7$, $f'(1)=2$, $f''(1)=16$, then $\left(f^{-1}\right)''(7)$ is equal to
\begin{qoptions}
\opt{A}{$-4$}
\opt{B}{$-\dfrac12$}
\opt{C}{$-1$}
\opt{D}{$-2$}
\end{qoptions}

\qn{12} $\displaystyle\lim_{t\to0}\dfrac{12^t-3^t-4^t+1}{\sqrt{2\cos t+7}-3}$ is equal to
\begin{qoptions}
\opt{A}{$\log_e12$}
\opt{B}{$\log_e3.\log_e4$}
\opt{C}{$\left(\log_e3.\log_e4\right)^{-1}$}
\opt{D}{$-6\log_e3.\log_e4$}
\end{qoptions}

\qn{13} If the function $f(x)=\begin{cases}2x & ;\ |x|\leqslant1\\ x^2+ax+b & ;\ |x|>1\end{cases}$ is continuous for all real $x$, then
\begin{qoptions}
\opt{A}{$a=2,\ b=-1$ and $f$ is differentiable for all $x$}
\opt{B}{$a=-2,\ b=1$ and $f$ is not differentiable at $x=-1,1$}
\opt{C}{$a=2,\ b=-1$ and $f$ is not differentiable at $x=-1,1$}
\opt{D}{$a=-2,\ b=-1$ and $f$ is not differentiable at $x=-1,1$}
\end{qoptions}

\qn{14} How many 6 letter words can be formed by using the letters of the word ``GENIUS'' such that each word neither start with `G' nor end with `S'.
\begin{qoptions}
\opt{A}{480}
\opt{B}{600}
\opt{C}{504}
\opt{D}{720}
\end{qoptions}

\qn{15} Let $\vec a,\vec b$ be two non collinear unit vectors and $\vec x=\vec a-\left(\vec a.\vec b\right)\vec b$ and $\vec y=\vec a\times\vec b$. Now consider the following statements :\\
\textbf{Statement-I :} $|\vec x|=|\vec y|$\\
\textbf{Statement-II :} $|\vec y|=|\vec x|+\left|\vec x.\vec b\right|$
\begin{qoptions}
\opt{A}{Statement-I is true but Statement-II is false}
\opt{B}{Both Statement-I and Statement-II are false}
\opt{C}{Both Statement-I and Statement-II are true}
\opt{D}{Statement-I is false but Statement-II is true}
\end{qoptions}

\qn{16} If $I_1=\displaystyle\int_0^{\pi/2}\dfrac{\sqrt{\cot x}}{\sqrt{\cot x}+\sqrt{\tan x}}\,dx$ \& $I_2=\displaystyle\int_{-\pi}^{\pi}\dfrac{\cos^2x}{1+a^x}\,dx$, then -
\begin{qoptions}
\opt{A}{$\dfrac1{I_1}-\dfrac1{I_2}=\dfrac3\pi$}
\opt{B}{$\dfrac2{I_1}-\dfrac1{I_2}=\dfrac5\pi$}
\opt{C}{$\dfrac2{I_1}+\dfrac1{I_2}=\dfrac8\pi$}
\opt{D}{$\dfrac1{I_1}+\dfrac1{I_2}=\dfrac6\pi$}
\end{qoptions}

\qn{17} If $I=\displaystyle\int_0^1\dfrac{dx}{1+x^{\pi/2}}$, then
\begin{qoptions}
\opt{A}{$I=\log_e2$}
\opt{B}{$I<\log_e2$}
\opt{C}{$I=\dfrac\pi4$}
\opt{D}{$\log_e2<I<\dfrac\pi4$}
\end{qoptions}

\qn{18} The smaller area (in sq. units) included between the curves $\sqrt x+\sqrt{|y|}=1$ and $|x|+|y|=1$ is
\begin{qoptions}
\opt{A}{$\dfrac13$}
\opt{B}{$\dfrac43$}
\opt{C}{$\dfrac23$}
\opt{D}{$\dfrac53$}
\end{qoptions}

\qn{19} Area bounded by $y=\sin x$ and $y-\dfrac12=\dfrac3{2\pi}\left(x-\dfrac\pi6\right)$ is equal to ($x>0$)
\begin{qoptions}
\opt{A}{$\dfrac{\sqrt3}2-\dfrac\pi4$}
\opt{B}{$\sqrt3-\dfrac\pi2$}
\opt{C}{$\dfrac{\sqrt3}2-\dfrac\pi6$}
\opt{D}{$2\sqrt3-\pi$}
\end{qoptions}

\qn{20} Diagonals of a quadrilateral are equal in length and they are perpendicular. Length of its diagonals is 6 units, consider four statements regarding this quadrilateral\\
I \quad This quadrilateral is a square\\
II \quad Area of this quadrilateral is 18 square units.\\
III \quad Diagonals bisect each other\\
IV \quad A diagonal divide the quadrilateral in two triangles of equal area.
\begin{qoptions}
\opt{A}{All 4 statements are correct}
\opt{B}{Exactly 3 statements are correct}
\opt{C}{Exactly 2 statements are correct}
\opt{D}{Exactly 1 statement is correct}
\end{qoptions}

\sect{Section-II}

\qn{1} The mean weight of 9 items is 15. If one more item is added to the series, the mean becomes 16. The value of 10$^{\text{th}}$ item is

\qn{2} If $A=\left\{x\,|\,x\in N\text{ and }x<6\tfrac14\right\}$ and $B=\left\{x\,|\,x\in N\text{ and }x^2\leqslant5\right\}$. Then the number of subsets of set $A\times(A\cap B)$ which contains 3 elements is

\qn{3} Let the equation of two ellipses be\\
$E_1:\dfrac{x^2}{3}+\dfrac{y^2}{2}=1$ and $E_2:\dfrac{x^2}{16}+\dfrac{y^2}{b^2}=1$ ($b<4$). If the product of their eccentricities is $\dfrac12$, then the length of the minor axis of ellipse $E_2$ is

\qn{4} Let $\{a_n\}$ and $\{b_n\}$ be two different sequences such that $a_{n+1}=2+a_n\ \forall\ n\in N$, $a_{10}=21$ and $b_n=\dfrac{1}{a_na_{n+2}}\ \forall\ n\in N$, the value of $15\displaystyle\sum_{r=1}^{\infty}b_r$ is

\qn{5} If $z_1$ and $z_2$ are non zero solutions of equation $z^2+z=i.\bar z$ where $i=\sqrt{-1}$, then $|z_1+z_2|$ is equal to

\qn{6} A variable chord of the hyperbola $\dfrac{x^2}{4}-\dfrac{y^2}{8}=1$, subtends a right angle at the centre of the hyperbola. If this chord touches a fixed circle which is concentric with the hyperbola and $r$ is radius of the circle, then $r^2$ is

\qn{7} If $f(x)=x^3+\log_2\left(x+\sqrt{x^2+1}\right)$ and $f(a)+f(b)\geq0$, then the least value of $(a+b)$ is

\qn{8} All the terms of an increasing geometric progression are positive integers and no two of them are equal. If product of its first five terms is P and $P=(12)^5$, then the common ratio of this G.P. is :

\qn{9} Let $z_1,z_2,z_3\in C$ such that $|z_1|=1$, $\left|z_2-\left(1+\sqrt3i\right)\right|=3$ and $\left|z_3-3-3\sqrt3i\right|=7$. If $|3z_1-2z_2-4z_3|$ is maximum then the value of $|z_3-z_1|+\left|\dfrac{z_3+3z_1}{z_2}\right|$ is equal to

\qn{10} The value of $\displaystyle\lim_{n\to\infty}\dfrac{\left(\displaystyle\lim_{x\to0^-}\sum_{r=1}^{2n+1}[x^r]\right)+n+3}{\ell n(1+n)^n-\ell n(n)^n}$ ; where $[.]$ denotes greatest integers functions

\end{multicols}
\end{document}
